\chapter{Divide and Conquer}

\begin{important}[Paradigm]
\textbf{Divide} the instance into smaller instances of the same problem,
\textbf{conquer} them recursively, and \textbf{combine} the sub-solutions.
Correctness is by (strong) induction on the size; runtime is from a recurrence.
\end{important}

\section{Mergesort and counting inversions}

\subsection{Mergesort}
\begin{lstlisting}[style=pseudo]
MergeSort(A[1..n]):
    if n <= 1 then return A
    L = MergeSort(A[1..n/2]);  R = MergeSort(A[n/2+1..n])
    return Merge(L, R)

Merge(L, R):              // L, R sorted
    i = 1; j = 1; out = []
    while i <= |L| and j <= |R| do
        if L[i] <= R[j] then append L[i] to out; i = i + 1
        else                 append R[j] to out; j = j + 1
    append the rest of L and R to out
    return out
\end{lstlisting}
\code{Merge} takes $\Theta(|L|+|R|)$, so $T(n)=2T(n/2)+\Theta(n)\in\Theta(n\log n)$.
Correctness of \code{Merge}: invariant ``\code{out} holds the $i+j-2$ smallest elements of $L\cup R$, in sorted order''.

\subsection{Counting inversions}
\begin{definition}
An \emph{inversion} in $A[1..n]$ is a pair $(i,j)$ with $i<j$ and $A[i]>A[j]$.
\end{definition}
The number of inversions measures how far $A$ is from sorted (between $0$ and $\binom n2$);
it is used, e.g., to compare two rankings. Brute force is $\Theta(n^2)$.

Divide and conquer: inversions are of three types: both indices in the left half,
both in the right half, or \emph{split} ($i$ in the left, $j$ in the right).
The first two are counted recursively. To count split inversions quickly, sort both halves
(which doesn't change the split count) and count during the merge:
when $R[j]$ is output before the remaining elements $L[i..|L|]$, it forms an inversion with
each of them, i.e.\ with $|L|-i+1$ elements.

\begin{lstlisting}[style=pseudo]
SortCount(A):                    // returns (sorted A, #inversions)
    if |A| <= 1 then return (A, 0)
    (L, a) = SortCount(left half);  (R, b) = SortCount(right half)
    i = 1; j = 1; c = 0; out = []
    while i <= |L| and j <= |R| do
        if L[i] <= R[j] then append L[i]; i = i + 1
        else append R[j]; c = c + (|L| - i + 1); j = j + 1
    append leftovers
    return (out, a + b + c)
\end{lstlisting}
Same recurrence as mergesort: $\Theta(n\log n)$.

\section{Integer multiplication: Karatsuba}
Multiplying two $n$-bit integers by the grade-school method takes $\Theta(n^2)$ bit operations.
Split $x=x_1 2^{n/2}+x_0$ and $y=y_1 2^{n/2}+y_0$. Then
\[
xy = x_1y_1 2^n + (x_1y_0+x_0y_1)2^{n/2}+x_0y_0.
\]
Four recursive multiplications of $n/2$-bit numbers plus $\Theta(n)$ for shifts/additions gives
$T(n)=4T(n/2)+\Theta(n)\in\Theta(n^2)$: no gain.

\textbf{Karatsuba's trick} (1960): the middle term is
\[
x_1y_0+x_0y_1=(x_1+x_0)(y_1+y_0)-x_1y_1-x_0y_0,
\]
so three multiplications suffice: $P_1=x_1y_1$, $P_2=x_0y_0$, $P_3=(x_1+x_0)(y_1+y_0)$, and
$xy=P_1 2^n+(P_3-P_1-P_2)2^{n/2}+P_2$.
\begin{lstlisting}[style=pseudo]
Karatsuba(x, y, n):
    if n <= threshold then return x * y        // grade school
    split x = (x1, x0), y = (y1, y0) at n/2 bits
    P1 = Karatsuba(x1, y1, n/2)
    P2 = Karatsuba(x0, y0, n/2)
    P3 = Karatsuba(x1 + x0, y1 + y0, n/2 + 1)
    return P1 * $2^n$ + (P3 - P1 - P2) * $2^{n/2}$ + P2
\end{lstlisting}
$T(n)=3T(n/2)+\Theta(n)$, so $T(n)\in\Theta(n^{\log_2 3})\approx\Theta(n^{1.585})$ by Master case 1.
(The extra bit in $x_1+x_0$ does not affect the asymptotics.)
Toom--Cook and FFT-based methods do better (down to $\bigO(n\log n)$, Harvey--van der Hoeven 2019).

\section{Matrix multiplication: Strassen}
Naive multiplication of $n\times n$ matrices is $\Theta(n^3)$. Splitting into $n/2\times n/2$ blocks,
\[
\begin{pmatrix}A&B\\C&D\end{pmatrix}\begin{pmatrix}E&F\\G&H\end{pmatrix}
=\begin{pmatrix}AE+BG & AF+BH\\ CE+DG & CF+DH\end{pmatrix}
\]
uses 8 block products: $T(n)=8T(n/2)+\Theta(n^2)=\Theta(n^3)$.
\textbf{Strassen} (1969) uses 7 products:
\begin{align*}
M_1&=(A+D)(E+H), & M_2&=(C+D)E, & M_3&=A(F-H), & M_4&=D(G-E),\\
M_5&=(A+B)H, & M_6&=(C-A)(E+F), & M_7&=(B-D)(G+H),
\end{align*}
and then
\begin{align*}
AE+BG&=M_1+M_4-M_5+M_7, & AF+BH&=M_3+M_5,\\
CE+DG&=M_2+M_4, & CF+DH&=M_1-M_2+M_3+M_6.
\end{align*}
(You are not expected to memorize these; check one entry, e.g.\ $M_3+M_5=AF-AH+AH+BH$.)
$T(n)=7T(n/2)+\Theta(n^2)\in\Theta(n^{\log_2 7})\approx\Theta(n^{2.81})$.
The point: \emph{reducing the number of recursive calls} changes the exponent; extra additions only affect the $f(n)$ term.

\section{Closest pair of points}
\begin{definition}
Given $n$ points in the plane, find a pair with minimum Euclidean distance.
\end{definition}
Brute force is $\Theta(n^2)$. In 1D, sort and check neighbours: $\Theta(n\log n)$.

\paragraph{Algorithm.} Presort the points by $x$ (array $P_x$) and by $y$ (array $P_y$).
\begin{enumerate}
  \item Split by the median $x$-coordinate into $Q$ (left) and $R$ (right), each of size $n/2$.
        Also split $P_y$ into $Q_y$, $R_y$ in linear time (preserving $y$-order).
  \item Recursively find $\delta_Q$, $\delta_R$; let $\delta=\min(\delta_Q,\delta_R)$.
  \item Any closer pair must be split, with both points within $\delta$ of the dividing line $x=x^*$.
        Let $S_y$ be the points of $P_y$ with $\abs{x-x^*}<\delta$, in $y$-order.
  \item For each point in $S_y$, compare it with the next 7 points in $S_y$. Return the best found.
\end{enumerate}

\begin{lemma}
If $p,q\in S_y$ and $d(p,q)<\delta$, then $p$ and $q$ are within 7 positions of each other in $S_y$.
\end{lemma}
\begin{proof}
Consider the $2\delta\times\delta$ rectangle above $p$ (from $p$'s $y$ up to $y+\delta$, spanning the strip).
Divide it into 8 squares of side $\delta/2$, four on each side of the line. Each square lies entirely on one side,
and any two points in a square are at distance at most $\delta/\sqrt2<\delta$, contradicting the definition of
$\delta$ (they would be on the same side). So each square holds at most one point, and the rectangle holds at most 8
points including $p$. Hence any $q$ above $p$ with $d(p,q)<\delta$ is among the next 7 entries.
\end{proof}

Runtime: presorting $\bigO(n\log n)$ once; then $T(n)=2T(n/2)+\Theta(n)\in\Theta(n\log n)$.
(Re-sorting by $y$ at each level would give $\Theta(n\log^2n)$.)

\section{Selection: median of medians}
\textbf{Selection problem}: given an unsorted array $A$ of $n$ distinct numbers and $k$, return the $k$-th smallest.
Sorting gives $\Theta(n\log n)$; we want $\Theta(n)$.

\textbf{Quickselect}: pick a pivot $p$, partition into $L=\{x<p\}$, $G=\{x>p\}$; if $|L|=k-1$ return $p$,
if $k\le|L|$ recurse on $L$, else recurse on $G$ with $k-|L|-1$.
With a random pivot the expected time is $\Theta(n)$, but the worst case is $\Theta(n^2)$.

\textbf{Median of medians} (Blum, Floyd, Pratt, Rivest, Tarjan 1973) chooses a pivot guaranteed to be good:
\begin{lstlisting}[style=pseudo]
Select(A, k):
    if |A| <= 5 then sort A and return A[k]
    divide A into $\lceil n/5\rceil$ groups of 5; find each group's median   // $\Theta(n)$
    M = list of these medians
    p = Select(M, $\lceil |M|/2 \rceil$)           // median of medians
    partition A around p into L, {p}, G
    if k == |L| + 1 then return p
    elif k <= |L| then return Select(L, k)
    else return Select(G, k - |L| - 1)
\end{lstlisting}

\begin{lemma}
$|L|\le \frac{7n}{10}+\bigO(1)$ and $|G|\le\frac{7n}{10}+\bigO(1)$.
\end{lemma}
\begin{proof}
About half of the $n/5$ groups, i.e.\ at least $\ceil{\frac12\ceil{n/5}}$ groups, have median $\ge p$.
Each such group (except possibly $p$'s own group and a final partial group) contributes 3 elements $\ge p$
(its median and the two larger ones). So at least $3(\frac{n}{10}-2)=\frac{3n}{10}-6$ elements are $\ge p$,
hence $|L|\le\frac{7n}{10}+6$. Symmetrically for $G$.
\end{proof}

So $T(n)\le T(n/5)+T(7n/10+6)+\Theta(n)$. Since $\frac15+\frac7{10}=\frac9{10}<1$, the theorem from Chapter 2 gives
$T(n)\in\Theta(n)$. (With groups of 3 the fractions sum to $1$ and we only get $\Theta(n\log n)$; groups of 7 also work.)

\section*{Summary}
\begin{itemize}
  \item Inversions: count split inversions during the merge; $\Theta(n\log n)$.
  \item Karatsuba: 3 half-size products, $\Theta(n^{\log_23})$. Strassen: 7 products, $\Theta(n^{\log_27})$.
  \item Closest pair: strip argument, each point checks $\le7$ neighbours in $y$-order; $\Theta(n\log n)$.
  \item Median of medians: pivot is guaranteed to discard $\approx 3n/10$ elements, giving $T(n)=T(n/5)+T(7n/10)+\Theta(n)=\Theta(n)$.
\end{itemize}
